\chapter{RLHF}
\label{ch:rlhf}

% Core question: How do you use RL to align a language model to human preferences?
%
% Length estimate: 20-25 pages
% Writing order: 4th (builds on Ch.13 preference data/RM + Ch.14 RL machinery)
%
% InstructGPT thread: MAIN CASE STUDY role. By end of this chapter, students have
% InstructGPT as a fully-developed mental model — but understood as historical
% foundation, not current best practice.

\textit{Placeholder --- full content follows the section outline below.}

\section{The Three-Stage Pipeline}
\label{sec:three-stage}

% InstructGPT as anchor case study: SFT -> RM -> PPO.
% How this chapter builds on Ch.11 (SFT) and Ch.13 (RM, preference data).
% Ch.14 provides the RL vocabulary used throughout.

\section{Reward Model in the RLHF Loop}
\label{sec:rm-in-loop}

% Reference to Ch.13 reward model training.
% RM's role during PPO: frozen scoring function.
% RM as imperfect proxy for human preferences.

\section{PPO for Language Models}
\label{sec:ppo-lm}

% Ch.14 PPO instantiated on LLM trajectories.
% Per-token vs per-sequence reward attribution.
% Concrete: InstructGPT hyperparameters, training details.
% Implementation: value head, rollout buffer, mini-batch updates.

\section{KL Constraint Deep Dive}
\label{sec:kl-deep-dive}

% Why KL to pi_ref is essential (building on Ch.14 §14.7).
% Empirical pathologies without KL constraint.
% InstructGPT ablations.
% KL coefficient scheduling.

\section{Reward Hacking and Mitigations}
\label{sec:reward-hacking}

% Over-optimization phenomenon (Gao et al. 2022 scaling law).
% Concrete examples: length exploitation, formatting tricks.
% Reward model ensembling.
% KL adaptive scheduling.

\section{RLHF Variants}
\label{sec:rlhf-variants}

% Constitutional AI (Anthropic): RLAIF with principles.
% Llama-2 RLHF recipe.
% Other production recipes.

\section{Industry Retreat from Full RLHF}
\label{sec:rlhf-retreat}

% 2024+ open-source preferring SFT + DPO (callback to Ch.13 §13.13).
% Frontier labs still using RLHF for nuanced alignment.
% Honest framing: RLHF is operationally expensive, DPO covers 80% of use cases.

\section{Why We Still Learn RLHF}
\label{sec:why-learn-rlhf}

% Conceptual foundation for Ch.16 reasoning RL.
% Required for frontier alignment work.
% Understanding the full pipeline gives judgment about when DPO is insufficient.
% Forward reference to Ch.16.


%----------------------------------------------------------------------
% End matter
%----------------------------------------------------------------------

% \subsection*{Summary}
% \subsection*{Pause-and-Verify Exercises}
% \subsection*{Key Readings}
% \subsection*{Key Equations}
